% LaTeX explanation for Convergence Analysis plots
% CEE 231 - Homework 7

\section*{Convergence Analysis: Torsional Stiffness}

The torsional stiffness $k_T$ for a rectangular cross-section is computed using an infinite series representation involving only \textit{odd} indices. In practice, this series must be truncated to a finite number of terms $N$. It is crucial to distinguish between $N$ (the number of terms) and $n$ (the index values): since only odd indices are included, $N$ terms correspond to $n = 1, 3, 5, \ldots, (2N-1)$. For example, $N=50$ terms means evaluating the series up to $n=99$.

Figure~\ref{fig:convergence} presents a convergence analysis to determine the required number of terms for accurate computation, using the parameters $\mu = 82.0~\text{kN/mm}^2$, $b = 10.0~\text{mm}$, and $c = 20.0~\text{mm}$.

\subsection*{Upper Panel: Torsional Stiffness Convergence}

The upper panel displays the cumulative torsional stiffness $k_T(N)$ as a function of the number of series terms $N$ (where $N$ terms corresponds to odd indices $n = 1, 3, 5, \ldots, 2N-1$). Several key observations can be made:

\begin{itemize}
    \item \textbf{Rapid initial convergence:} The function exhibits steep growth for small values of $N$ (approximately $N < 50$, i.e., $n < 99$), indicating that the first few terms contribute the majority of the total stiffness. This is characteristic of series with terms decaying as $n^{-4}$, where higher-order terms have progressively diminishing contributions.
    
    \item \textbf{Asymptotic behavior:} Beyond approximately $N = 100$ (i.e., $n = 199$), the curve becomes nearly horizontal, asymptotically approaching the converged value of $k_T \approx 375.0~\text{kN·mm}^2$ (indicated by the dashed red line). This suggests that including terms beyond $N = 100$ provides negligible improvement in accuracy.
    
    \item \textbf{Practical truncation:} For engineering applications requiring typical accuracy ($\sim$0.1\%), using $N = 50$ terms (up to $n=99$) yields $k_T \approx 374.95~\text{kN·mm}^2$, which differs from the converged value by less than 0.01\%. This validates the common practice of using $N = 50$--$100$ terms (corresponding to maximum $n$ values of $99$--$199$) for such calculations.
\end{itemize}

\subsection*{Lower Panel: Relative Change Analysis}

The lower panel presents the relative change between consecutive terms, defined as:
\begin{equation}
    \text{Relative Change} = \left| \frac{k_T(N) - k_T(N-1)}{k_T(N)} \right| \times 100\%
\end{equation}

This metric quantifies the incremental contribution of each additional term. Key observations include:

\begin{itemize}
    \item \textbf{Logarithmic decay:} The relative change decreases exponentially with increasing $N$, appearing as a nearly linear trend on the semi-logarithmic plot. This demonstrates the series' rapid convergence properties.
    
    \item \textbf{Convergence rate:} At $N = 50$ (corresponding to $n = 99$), the relative change is approximately $10^{-4}\%$, indicating that the 50th term (with index $n=99$) changes the result by only one part in a million. By $N = 1000$ ($n = 1999$), the relative change drops below $10^{-10}\%$, effectively reaching machine precision limits.
    
    \item \textbf{Practical implications:} The steep initial decline (for $N < 20$, i.e., $n < 39$) confirms that a modest number of terms suffices for engineering accuracy. The continued monotonic decrease validates the series' mathematical convergence and numerical stability.
\end{itemize}

\subsection*{Conclusion}

The convergence analysis demonstrates that the infinite series representation of torsional stiffness converges rapidly and monotonically. For the given geometry and material properties, using $N = 50$ terms (evaluating odd indices up to $n = 99$) provides sufficient accuracy for practical engineering calculations, yielding a torsional stiffness of $k_T \approx 374.95~\text{kN·mm}^2$. The exponential decay of the relative change metric confirms that subsequent terms contribute negligibly, making higher truncation values unnecessary for most applications.

This behavior is consistent with the theoretical expectation that terms decaying as $O(n^{-4})$ in a Fourier-type series converge rapidly, requiring only a modest number of terms to achieve engineering precision. The distinction between $N$ (number of terms) and $n_{\text{max}} = 2N-1$ (maximum index value) is important for correctly implementing and interpreting the series. The analysis validates the numerical implementation and provides confidence in the computed torsional stiffness value.

